\section*{Chapitre \ref{ch:b3:origin-of-elements} --- D'où viennent les éléments}
% Values from the Book 5 scratch calculator calc/ch30.py.

\begin{solution}{exo:b3:origin-of-elements:1}
(a)~$\mathrm{Si/Fe} = 10^{7.51 - 7.46} = 10^{0.05} = 1.12$ et
$\mathrm{O/Fe} = 10^{8.69 - 7.46} = 10^{1.23} = 17.0$~: environ un atome de silicium
et dix-sept atomes d'oxygène pour un atome de fer. (b)~$\mathrm{Fe/Au} =
10^{7.46 - 0.91} = 10^{6.55} = \num{3.55e6}$ et $\mathrm{Pb/Au} =
10^{1.95 - 0.91} = 10^{1.04} = 11.0$~: au-delà du pic du fer, les abondances
chutent de plusieurs millions, mais le plomb, qui forme à lui seul un pic,
est onze fois plus abondant que l'or.
\end{solution}

\begin{solution}{exo:b3:origin-of-elements:2}
(a)~$10^{7.46 - 12} = \num{2.884e-5}$, soit \num{28.8} atomes de fer par
million d'atomes d'hydrogène. (b)~Cent fois moins~: \num{0.288} par
million, soit un atome de fer pour $1/\num{2.884e-7} = \num{3.47e6}$ atomes
d'hydrogène. (c)~$[\mathrm{Fe/H}] = \log_{10}(\num{e-7}/\num{2.884e-5}) =
-2.46$.
\end{solution}

\begin{solution}{exo:b3:origin-of-elements:3}
$^{124}\mathrm{Xe}$ et $^{126}\mathrm{Xe}$ se situent du côté riche en protons,
hors de la voie s, et sont protégés par les isotopes stables $^{124}\mathrm{Te}$ et
$^{126}\mathrm{Te}$~: \omterm{def:b3:origin-of-elements:p-process}{noyaux p}. $^{128}\mathrm{Xe}$ et $^{130}\mathrm{Xe}$
sont sur la voie s mais protégés du \omterm{def:b3:origin-of-elements:s-process}{processus r} par les isotopes stables
$^{128}\mathrm{Te}$ et $^{130}\mathrm{Te}$~: s seul. $^{129}\mathrm{Xe}$,
$^{131}\mathrm{Xe}$ et $^{132}\mathrm{Xe}$ sont sur la voie s et n'ont aucun isobare
stable de $Z$ inférieur (l'iode n'en a qu'un, à 127)~: les deux processus.
$^{134}\mathrm{Xe}$ et $^{136}\mathrm{Xe}$ se situent au-delà de
$^{133}\mathrm{Xe}$, là où la voie s quitte le xénon, et ne sont pas
protégés~: r seul.
\end{solution}

\begin{solution}{exo:b3:origin-of-elements:4}
(a)~$t = d/c = 40\times\qty{3.086e22}{m}/\qty{2.998e8}{m.s^{-1}} =
\qty{4.12e15}{s}$, soit environ \num{130} millions d'années. (b)~$\Delta v/c \le
1.7/\num{4.12e15} = \num{4.1e-16}$. Il s'agit d'une borne supérieure, car une
partie ou la totalité du retard peut provenir de l'émission des rayons gamma
après la fusion.
\end{solution}

\begin{solution}{exo:b3:origin-of-elements:5}
(a)~$295\times0.1067 = \qty{31.48}{mb}$, $832\times0.0373 =
\qty{31.03}{mb}$, $155\times0.1984 = \qty{30.75}{mb}$. (b)~À partir de
$^{122}\mathrm{Te}$~: $N_{123} = 31.48/832 = \num{0.03783}$, \qty{1.4}{\percent}
au-dessus de la valeur mesurée, et $N_{124} = 31.48/155 = \num{0.2031}$,
\qty{2.4}{\percent} au-dessus. Des abondances qui diffèrent d'un facteur cinq
sont prédites à quelques pour cent près par les seules sections efficaces.
(c)~$^{128}\mathrm{Te}$ est r seul~: le flux s quitte le tellure
à $^{127}\mathrm{Te}$, qui se désintègre, de sorte qu'aucun flux s stationnaire
traverse $^{128}\mathrm{Te}$ et son abondance a été fixée par le \omterm{def:b3:origin-of-elements:s-process}{processus r}.
\end{solution}

\begin{solution}{exo:b3:origin-of-elements:6}
(a)~$\tau_0\langle\sigma\rangle = 0.306\times155 = 47.4$, donc
$f = 1/(1 + 1/47.4) = \num{0.979}$ pour $^{124}\mathrm{Te}$~; pour
$^{138}\mathrm{Ba}$, $0.306\times4.0 = 1.224$ et $f = 1/(1 + 0.817) =
\num{0.550}$~: le noyau magique ne laisse passer que \qty{55}{\percent} du
flux. (b)~Les six facteurs valent \num{0.959}, \num{0.550}, \num{0.908},
\num{0.771}, \num{0.972} et \num{0.915}~; leur produit vaut \num{0.328}~:
$\sigma N$ diminue d'un facteur \num{3.05}. (c)~La chute mesurée vaut
$27.6/7.72 = 3.57$. Le modèle reproduit l'ampleur de la marche à
\qty{15}{\percent} près, presque entièrement due à $^{138}\mathrm{Ba}$ et
$^{140}\mathrm{Ce}$, les deux noyaux de $N = 82$.
\end{solution}

\begin{solution}{exo:b3:origin-of-elements:7}
(a)~$kT = 8\times10^3\times\qty{1.602e-19}{J} = \qty{1.282e-15}{J}$, et
$v_T = \sqrt{2\times\num{1.282e-15}/\num{1.6749e-27}} =
\qty{1.24e6}{m.s^{-1}}$. (b)~Avec $\langle\sigma\rangle =
\qty{155e-31}{m^2}$, $\lambda = n_\mathrm{n}\langle\sigma\rangle v_T$
vaut \qty{1.92e-11}{s^{-1}} pour $n_\mathrm{n} = \qty{e12}{m^{-3}}$
($10^6$ par centimètre cube)~: un temps moyen de \qty{5.2e10}{s}, soit environ
\num{1650} ans~; puis \num{16.5} ans à $10^8$ et \num{0.165}~an,
deux mois, à \qty{e10}{cm^{-3}}. (c)~$^{127}\mathrm{Te}$ se désintègre en
heures, bien plus vite que chacun de ces temps de capture~: il se désintègre
toujours, et la voie ne se divise pas à cet endroit. $^{134}\mathrm{Cs}$ vit environ deux
ans, entre les temps de capture à $10^8$ et \qty{e10}{cm^{-3}}~: dans les
fortes densités des flashs d'hélium, une bonne partie de ce noyau capture
un neutron avant de se désintégrer. C'est un point de branchement, et les abondances
de part et d'autre permettent de mesurer la densité neutronique.
\end{solution}

\begin{solution}{exo:b3:origin-of-elements:8}
(a)~$10^{6.30 - 3.14} = 1445$ fois plus rare que le calcium et
$10^{4.97 - 3.14} = 68$ fois plus rare que le titane. (b)~$10^{8.46 - 1.38}
= \num{1.2e7}$. (c)~Le scandium présente l'effet pair-impair~: les noyaux
à nombre pair de protons sont plus fortement liés et sont construits par la
capture de noyaux d'hélium pendant les phases de combustion des étoiles massives,
qui les favorisent. Le béryllium illustre le creux du lithium, du béryllium et du bore~:
ces noyaux fragiles sont détruits aux températures de l'intérieur stellaire,
et la combustion de l'hélium passe directement de l'hélium au carbone,
sans les produire.
\end{solution}

\begin{solution}{exo:b3:origin-of-elements:9}
(a)~Notons $\alpha$ les \omterm{def:b3:origin-of-elements:alpha}{éléments alpha} et $\mathrm{Fe}_\mathrm{cc}$,
$\mathrm{Fe}_\mathrm{Ia}$ le fer issu des deux types de supernovae. Sur
le palier, $\alpha/\mathrm{Fe}_\mathrm{cc} = 10^{P}$ fois le rapport
solaire. À $[\mathrm{Fe/H}] = 0$, $\alpha/(\mathrm{Fe}_\mathrm{cc} +
\mathrm{Fe}_\mathrm{Ia})$ est solaire. En divisant,
$\mathrm{Fe}_\mathrm{cc}/(\mathrm{Fe}_\mathrm{cc} + \mathrm{Fe}_\mathrm{Ia})
= 10^{-P}$, et la part des supernovae de type~Ia vaut $1 - 10^{-P} = 1 - 10^{-0.4} =
\num{0.602}$. (b)~$10^{-P} = 1/3$ donne $P = \log_{10}3 = 0.48$.
(c)~Les supernovae de type~Ia commencent après le même délai dans les deux
galaxies. Lorsqu'elles ont débuté dans la galaxie naine, son gaz n'avait
atteint que $[\mathrm{Fe/H}] = -2$, contre $-1$ dans la Galaxie~: elle avait
transformé son gaz en étoiles dix fois plus lentement.
\end{solution}

\begin{solution}{exo:b3:origin-of-elements:10}
(a)~Pour le zirconium,
\[
\frac{(X/\mathrm{Mg})_\mathrm{Earth}}{(X/\mathrm{Mg})_\mathrm{CI}} =
\frac{10.5/228000}{3.96/95870} = 1.11~;
\]
de même $(2670/228000)/(5010/95870) = 0.224$ pour le sodium,
\num{0.190} pour le potassium et \num{0.118} pour le rubidium. (b)~Le zirconium, le plus réfractaire, se place
près de sa proportion chondritique. Les trois métaux alcalins
modérément volatils, qui se condensent à \qtyrange{800}{1006}{K}, sont appauvris par
un facteur de cinq à neuf~; le rubidium, le plus volatil, l'est davantage. La
Terre silicatée suit la \omterm{def:b3:origin-of-elements:volatility}{tendance de volatilité}. (c)~Le magnésium est un élément majeur,
réfractaire qui se condense avec l'essentiel de la roche, et il est
lithophile~: il est resté dans le manteau lors de la formation du noyau. Un élément sidérophile
tel que le fer ou le nickel mêlerait la perte vers le noyau à la
perte par volatilité.
\end{solution}

\begin{solution}{exo:b3:origin-of-elements:11}
(a)~Lorsqu'une masse $\mathrm{d}s$ est bloquée dans les étoiles, le vent retire
$\lambda\,\mathrm{d}s$ de gaz de \omterm{def:b3:origin-of-elements:metallicity}{métallicité} $Z$~:
$\mathrm{d}M_\mathrm{g} = -(1 + \lambda)\,\mathrm{d}s$ et
$\mathrm{d}(ZM_\mathrm{g}) = y\,\mathrm{d}s - Z\,\mathrm{d}s -
\lambda Z\,\mathrm{d}s$. En développant la seconde et en soustrayant
$Z\,\mathrm{d}M_\mathrm{g}$ donne $M_\mathrm{g}\,\mathrm{d}Z =
y\,\mathrm{d}s = -\frac{y}{1 + \lambda}\,\mathrm{d}M_\mathrm{g}$, ce qui
s'intègre pour donner le résultat. (b)~$\ln(1/0.29) = 1.238$~: $Z = 0.0155$ pour
$\lambda = 0$, et deux fois moins, $0.0077$, pour $\lambda = 1$. (c)~Un vent
abaisse le rendement effectif d'un facteur $1 + \lambda$. Dans une galaxie de
faible masse, les supernovae chassent plus facilement le gaz, $\lambda$ est grand,
et les métaux sont perdus avec le gaz.
\end{solution}

\begin{solution}{exo:b3:origin-of-elements:12}
(a)~$N_\mathrm{Li}/N_\mathrm{H} = 10^{0.96 - 12} = \num{9.1e-12}$ dans le
Soleil et $10^{3.25 - 12} = \num{1.78e-9}$ dans les chondrites CI. Le Soleil
conserve $10^{-2.29} = \num{0.0051}$ de son lithium~: il en a détruit
\qty{99.5}{\percent}. (b)~$\num{1.78e-9}/\num{1.45e-10} = 12.3$ et
$\num{1.78e-9}/\num{4.72e-10} = 3.8$. (c)~Le Système solaire contient près de
quatre fois plus de lithium que n'en a produit le Big Bang~: l'essentiel de son lithium a été
produit plus tard, par un autre processus. Les étoiles les plus anciennes en contiennent trois fois
moins que n'en a produit le Big Bang~: soit elles en ont détruit une partie, comme
le Soleil l'a fait, soit la prédiction omet quelque chose~; les deux comparaisons
ne permettent pas de trancher.
\end{solution}

\begin{solution}{pb:b3:origin-of-elements:1}

\textbf{Partie I.}

\textbf{1.} $10^{12 - 7.46} = 10^{4.54} = \num{34700}$.

\textbf{2.} $N_\mathrm{He}/N_\mathrm{H} = 10^{10.914 - 12} = \num{0.0820}$, d'où
$4N_\mathrm{He}/N_\mathrm{H} = \num{0.328}$, contre $Y/X =
0.2423/0.7438 = \num{0.326}$. Les deux valeurs s'accordent à \qty{0.7}{\percent}~; le
petit écart tient à ce qu'un atome d'hélium pèse un peu moins que quatre
atomes d'hydrogène.

\textbf{3.} $\log_{10}\bigl((1/350000)/\num{2.884e-5}\bigr) = -5.544 + 4.540 =
-1.00$.

\textbf{4.} $[\mathrm{O/H}] = -1.00 + 0.40 = -0.60$, d'où
$\log\varepsilon(\mathrm{O}) = 8.69 - 0.60 = 8.09$. L'étoile possède
$10^{8.09 - 6.46} = 42.7$ atomes d'oxygène par atome de fer, contre
$10^{1.23} = 17.0$ dans le Soleil~: $10^{0.4} = 2.5$ fois plus.

\textbf{5.} $Z \approx 0.0139\times10^{-0.60} = \num{0.0035}$, soit un quart de
celle du Soleil.

\textbf{6.} Le fer possède de très nombreuses raies dans le spectre d'une étoile froide, alors que
l'oxygène en a peu, et elles sont faibles~: le fer se mesure bien plus facilement et
avec plus de précision, et sert de référence pour les métaux.

\textbf{Partie II.}

\textbf{7.} Chacun se place derrière un isobare stable de $Z$ inférieur~:
$^{122}\mathrm{Sn}$, $^{123}\mathrm{Sb}$ et $^{124}\mathrm{Sn}$ arrêtent
les chaînes de désintégration du \omterm{def:b3:origin-of-elements:s-process}{processus r} avant qu'elles n'atteignent le tellure. Les trois se situent
sur la voie s, qui traverse le tellure de $^{122}\mathrm{Te}$ à
$^{126}\mathrm{Te}$.

\textbf{8.} \num{31.48}, \num{31.03} et \qty{30.75}{mb}~; moyenne
\qty{31.09}{mb}.

\textbf{9.} La moyenne des deux autres vaut $(31.48 + 30.75)/2 = \qty{31.11}{mb}$,
d'où $N_{123} = 31.11/832 = \num{0.03740}$, seulement \qty{0.3}{\percent}
au-dessus de la valeur mesurée \num{0.0373}.

\textbf{10.} $1.329\times44.4 = \qty{59.0}{mb}$, presque le double du palier.
$^{128}\mathrm{Te}$ n'est produit que par le \omterm{def:b3:origin-of-elements:s-process}{processus r}~: la voie s quitte le tellure au
$^{127}\mathrm{Te}$, son abondance a donc été produite par le \omterm{def:b3:origin-of-elements:s-process}{processus r} et
ne doit rien au flux s.

\textbf{11.} $1 - 7.43/8.53 = \num{0.13}$~: \qty{13}{\percent} du flux est passé
en contournant $^{148}\mathrm{Sm}$, par captures sur les
noyaux instables $^{147}\mathrm{Nd}$, $^{147}\mathrm{Pm}$ et $^{148}\mathrm{Pm}$ avant
de se désintégrer en samarium.

\textbf{12.} Avec une couche fermée de 82 neutrons, ces noyaux lient un neutron de plus
faiblement et ont des sections efficaces de capture très petites (\qty{4.0}{mb}
pour $^{138}\mathrm{Ba}$). Le flux s'accumule sur eux, et seule une partie
en repart~: ce sont des goulots d'étranglement.

\textbf{Partie III.}

\textbf{13.} $\mu_1 = 13.7/(13.7 + 33.4) = 13.7/47.1 = \num{0.291}$.

\textbf{14.} $\ln(1/0.291) = \num{1.235}$, d'où $y = 0.0154/1.235 = \num{0.0125}$,
\num{0.81} fois $Z_\odot$.

\textbf{15.} À $\mu = 0.9$, $\ln(1/0.9) = 0.1054$, donc $Z = 0.0125\times0.1054 =
\num{0.00131}$. Comme $0.00131/0.0154 = 0.0851$, $[\mathrm{Fe/H}] =
-1.07$. À $\mu = 0.5$, $Z = 0.0125\times0.693 = \num{0.00864}$ et
$[\mathrm{Fe/H}] = -0.25$.

\textbf{16.} $Z = 0.00154$, donc $\mu = \mathrm{e}^{-0.00154/0.01247} =
\mathrm{e}^{-0.1235} = \num{0.884}$.

\textbf{17.} $1 - 0.884 = \num{0.116}$~: \qty{12}{\percent} de la masse.

\textbf{18.} $y = 0.0154/\ln 10 = \num{0.0067}$, soit environ la moitié.

\textbf{Partie IV.}

\textbf{19.} La masse des étoiles de \omterm{def:b3:origin-of-elements:metallicity}{métallicité} inférieure à $Z$ vaut
$M(1 - \mathrm{e}^{-Z/y})$, donc par unité de $Z$ elle vaut
$(M/y)\,\mathrm{e}^{-Z/y}$. Comme $Z = Z_\odot10^{[\mathrm{Fe/H}]}$,
$\mathrm{d}Z = Z\ln 10\,\mathrm{d}[\mathrm{Fe/H}]$, et par unité de
$[\mathrm{Fe/H}]$ la distribution est proportionnelle à
$Z\mathrm{e}^{-Z/y}$. Sa dérivée,
$(1 - Z/y)\mathrm{e}^{-Z/y}$, s'annule en $Z = y$~: le maximum se situe en
$[\mathrm{Fe/H}] = \log_{10}0.810 = -0.09$.

\textbf{20.} $\num{0.116}/(1 - 0.291) = 0.1162/0.7091 = \num{0.164}$.

\textbf{21.} $Z/y = 1.235\times10^{-0.5} = 0.3905$, $1 - \mathrm{e}^{-0.3905} =
0.3233$, et $0.3233/0.7091 = \num{0.456}$.

\textbf{22.} $6670\times0.1638 = 1093$, contre 26~: quarante fois trop.

\textbf{23.} $26/902 = \qty{2.9}{\percent}$, soit encore plus de cinq fois moins que
dans la boîte. Les étoiles du modèle sont les plus durables, celles d'au plus
environ une masse solaire, qui ont toutes survécu depuis leur formation~; parmi
les étoiles F plus massives, les plus anciennes, qui sont aussi les plus pauvres
en métaux, ont disparu, et l'échantillon complet les sous-estime.

\textbf{24.} $(1 - 0.1^{0.1})/(1 - 0.1) = 0.2057/0.9 = \num{0.23}$~: pire.
Augmenter la fraction de gaz n'aide pas non plus~: lorsque $\mu_1$ tend vers 1,
la fraction tend vers $0.1$, soit encore plus de trois fois les
\qty{2.9}{\percent} observés. Aucune fraction de gaz ne sauve la boîte fermée.

\textbf{25.} La boîte fermée prédit que
\textbf{\qty[text-series-to-math=true]{16}{\percent}} des étoiles durables
du voisinage solaire ont $[\mathrm{Fe/H}] < -1$, contre \qty{3}{\percent}
au plus observés. Le disque n'était pas une boîte fermée~: il
s'est accru par un lent apport de gaz pauvre en métaux, de sorte que peu de gaz
est resté longtemps à basse \omterm{def:b3:origin-of-elements:metallicity}{métallicité} pour former des étoiles.
\end{solution}
